\section{Related Work and Limitations}
\label{sec:related}

This work builds directly on the backward-error methodology of
\citet{rasquinha2023metric}, extending it to block-scaled formats and
heavy-tailed operands. Two recent works approach FP4 microscaling from other
directions. \citet{fasoli2026finer} also measures block-size effects on model
correctness, but attributes degradation to \emph{narrow} tensor distributions ---
the opposite regime from the heavy tails swept here.
\citet{egiazarian2026bridging} analyses MXFP4 and NVFP4 error under a Laplace
operand model scored by per-element MSE, differing from this work in both metric
and operand model.

Six scope limits qualify these results. (i) The three per-tensor reference
formats --- FP8-E4M3, FP8-E5M2 and INT8 --- are in the frozen grid but never
executed, leaving MXFP4 and NVFP4 without a classical per-tensor baseline.
(ii) The random Hadamard transform is applied per block: a global, whole-row
transform would spread energy over $\sqrt{n}$, tying its effect to the same $n$
whose scaling law is under study, so global RHT is not evaluated here.
(iii) Quantization is CPU emulation in float64, not real quantized hardware.
(iv) Only E2M1 elements are studied. (v) Elements sharing a block share a
scale, so their errors are correlated rather than independent --- the leading
candidate behind the $u_{\mathrm{eff}}$ residual, consistent with a
survival-fraction gap whose rank order against the $u_{\mathrm{eff}}$ block-size
excess is perfect across all eight $\nu$ (Spearman 1.0), but not isolated by a
dedicated ablation. (vi) $c = 1$ is fixed a priori for
anti-circularity, so what is characterized is the bound's functional \emph{form},
not a calibrated absolute value.
